% =============================================================================
% Abkürzungsverzeichnis
% -----------------------------------------------------------------------------
% Nur FACHSPEZIFISCHE Abkürzungen aufnehmen. Triviale wie "z.B." oder "usw."
% gehören NICHT hierher. Alphabetisch sortiert halten.
%
% Im Text nutzt du \ac{KUERZEL}: Bei der ersten Verwendung wird der Begriff
% automatisch ausgeschrieben (inkl. Kürzel), danach erscheint nur das Kürzel.
% Dank [printonlyused] erscheinen hier NUR tatsächlich verwendete Abkürzungen.
% =============================================================================
\section*{Abkürzungsverzeichnis}

\begin{acronym}[DSGVO]  % längstes Kürzel als Referenzbreite für die Einrückung
  \acro{BSI}{Bundesamt für Sicherheit in der Informationstechnik}
  \acro{DSGVO}{Datenschutz-Grundverordnung}
  \acro{KI}{Künstliche Intelligenz}
  \acro{KMU}{kleine und mittlere Unternehmen}
  \acro{LLM}{Large Language Model}
\end{acronym}
